\chapter{Interpretation and Its Limits}
\label{comp:ch14}

This chapter covers \MT{ch:philosophy} together with the proof appendices that carry its four-cycle
results, from \MT{sec:phil-pairing-selection-proof} to \MT{sec:phil-early-parent}. The technical
chapter is short, but it is where the program's interpretive claims meet its most delicate formation
theorems, and eight of the propositions treated here are stated, not only proved, in those appendices.

\section{The question and the answer}
\label{comp:ch14-question}

The technical chapter asks what the preceding chapters commit to, and which readings of them an
observation or a proof could count against. Its answer has two layers that should be kept apart. The
first is interpretive and proves almost nothing: it sorts readings into declared rules of reading
(definitions), commitments the mathematics uses (standing hypotheses), and unsettled questions (open
problems). Its verdicts are that the architecture is local-first rather than background-free, that the
frame-change invariants are compatible with several structural-realist positions and select none, that
environment messages create no agents, and that nothing proved before the formation chapter
distinguishes the program from an ordinary hierarchical model with learned priors, whereas formation by
the inference could. The second layer is mathematical and substantial: a line of results on the
homogeneous four-cycle of \MT{def:rg-formation-cycle-experiment} asking whether learning, rather than a
random seed or a supplied hierarchy of strengths, decides which collective forms.

The answer of the second layer fits in one paragraph. On the four-cycle, formation in the technical
sense, meaning that a block enters coherence with a probability above its stationary value, happens
from every Gaussian initial law with a probability bounded below, but that probability is supplied by
the preparation's random initial assignment and source learning plays no role in it
(\MT{prop:phil-seeded-contrast-formation}). Law-level selection does happen when learning is slow: at a
small fixed temperature, as the learning mobility tends to zero, the learned source prior picks one of
the two pairings of the cycle for almost every Gaussian initial mean, while the arm with a frozen prior
ends symmetric (\MT{conj:phil-contrast-formation}, a theorem despite its label). That selection
amplifies an asymmetry already present in the initial law, it appears only at late times, and at those
times the selected pair's coherence is a property of equilibrium, not a formation event
(\MT{prop:phil-selected-pair-persistence}); on the early horizon where formation happens, slow learning
leaves the law of the first formed block unchanged (\MT{prop:phil-early-block-null}). Early formation
without law-level selection, late law-level selection without formation: the two halves do not meet,
and making them meet is the open problem \MT{open:rg-formation-adaptive-family}. The remaining results
concern the parent. The fixed parent channel of the four-cycle can be activated with separately
vanishing conditional and dynamical errors, both after the late selection and on the early seeded
episode, provided the activation time waits for the discarded modes to thermalize; activation at the
moment the coherence criterion is first met provably fails.

In the chain of the theory this chapter follows the formation chapter (\cref{comp:ch13}), whose
finite-depth theorem forms collectives through a supplied hierarchy of interaction strengths and whose
endogenous four-cycle witness forms one pairing from a symmetric preparation without learning. The
present chapter asks whether learning can do the choosing. \Cref{comp:tab-ch14-line} maps the
statements by time horizon and observable, the organizing principle the technical text uses but never
tabulates.

\begin{sciencetable}{The four-cycle continuation line. ``Early'' is the original episode horizon, of order $\log(1/\epsilon)$ at temperature $\mathrm T=\epsilon^8$; ``late'' is the selection window $t\sim\epsilon_{\mathrm L}^{-1}\log(1/\epsilon_{\mathrm L})$ at fixed or slowly vanishing temperature.}
\label{comp:tab-ch14-line}
\small
\begin{tabularx}{\textwidth}{>{\raggedright\arraybackslash}p{0.25\textwidth}>{\raggedright\arraybackslash}p{0.17\textwidth}>{\raggedright\arraybackslash}X>{\raggedright\arraybackslash}p{0.2\textwidth}}
\toprule
\textbf{Statement} & \textbf{Horizon} & \textbf{What is compared or bounded} & \textbf{Role of learning}\\
\midrule
\MT{prop:phil-seeded-contrast-formation} & early & probability of a contrast (entrance) event & idle \\
\MT{conj:phil-contrast-formation} & late, fixed $\mathrm T$ & one-time assignment law, learned versus frozen & amplifies the initial law \\
\MT{cor:phil-joint-selection} & late, $\mathrm T\to0$ jointly & same, along implicit mobility thresholds & same \\
\MT{prop:phil-temperature-tracking} & late & distance of the law from its instantaneous equilibrium & adiabatic following only \\
\MT{prop:phil-selected-pair-persistence} & late window $\mathrm T^{-1/4}$ & path law and coherent partition & fixes the label only \\
\MT{prop:phil-early-block-null} & early & conditional law of the first formed block & none if $\epsilon_{\mathrm L}=o(\mathrm T^{1/4})$ \\
\MT{prop:phil-actual-conditioned-parent} & late activation & discarded conditional information & parent inherits rows \\
\MT{prop:phil-actual-parent-dynamics} & late window $\mathrm T^{-1/4}$ & retained path law, parent versus fine & parent learns from its own law \\
\MT{prop:phil-early-seeded-parent} & early activation & discarded information, both directions & idle for $\epsilon_{\mathrm L}\in[0,1]$ \\
\MT{prop:phil-early-first-completion} & first completion & discarded information diverges & idle \\
\MT{prop:phil-early-parent-dynamics} & windows $\ll\mathrm T^{-1}$ & retained path law, parent versus fine & idle \\
\bottomrule
\end{tabularx}
\end{sciencetable}

\section{The picture}
\label{comp:ch14-picture}

\subsection{Scaffolding, idle posits, and what survives a change of frame}

Each agent carries two law points at the base point, a belief law and a model law, which the cross maps
do not identify (\MTeq{eq:phil-fiber-measures}). The base, the principal bundle, and the measurable
spaces are coordination scaffolding: they make the agents' interfaces comparable, but no inference law
constructs them and no law is placed on them. That is the meaning of local-first and not
background-free, and a physicist will recognize lattice gauge theory, where the lattice and the group
are fixed background and only site fields and links fluctuate. The verdict describes the present
architecture; it would change if edges became structural coordinates updated by inference
(\MT{open:obs-edge-creation}) or links were recovered from the records of a target population. Because the scaffolding
is not an inferential target, its lack of a law decides neither a Neo-Kantian reading (frames as
revisable constitutive conditions of meaning) nor a noumenal one (the bundle as structure reached only
through representations).

The chapter's one new definition is a parsimony rule: a posit is idle when changing it alters the law
of no record and of no statistic of a record in the model, and the idle-wheel criterion
(\MT{def:phil-idle-wheel}) removes idle posits. It plays the role gauge redundancy plays in physics, but
it is a declared rule of reading, close in spirit to constructive empiricism without being derived from
it, and it is relative to the records a model declares. What survives it is structural data,
\[
\mathcal I_{\mathrm{struct}}
=\Bigl\{
\text{ELBO and relative-entropy values},\
[\mathrm{Hol}^b(p)]_{\rm conj},\
[\mathrm{Hol}^m(p)]_{\rm conj},\
\bigl[(\mathring\Theta^b_e,\mathring\Theta^m_e)_{e}\bigr]
\Bigr\}
\]
(\MTeq{eq:phil-invariants}): scalars invariant because the kernels are covariant
(\MT{prop:gen-product-evidence-invariance}), loop classes in both channels, invariant because a passive
re-expression conjugates a loop product, and the per-edge twists (each link expressed in its endpoints'
own frames) modulo one simultaneous conjugation, which is finer than the loop classes because it keeps
how the twists of different edges sit relative to one another. For abelian links, as on the four-cycle,
the loop class is the loop sum $\sum_i\omega_i$. On a graph-flat family the loop classes are trivial
(\MT{prop:geo-trivializing-criterion}) and give a noumenal reading nothing to do; a nontrivial loop class
is a referent inside the interaction complex, which a declared record at $c_*$ can read
(\MT{thm:cg-loop-operational-holonomy}). The technical text finds these invariants compatible with
epistemic, moderate, and ontic structural realism and gives no argument from theory change, since they
are invariants of frame change inside one theory, whereas Worrall's structure \citep{Worrall1989}
survives a change of theory.

\begin{analogy}{exact}
The loop classes $[\mathrm{Hol}^x(p)]_{\rm conj}$ are the gauge-invariant content of a closed loop in
lattice gauge theory, the data of a Wilson loop before the trace. On the interaction graph this is an
identity of structure: links are group elements on edges, a passive re-expression acts as
$\Theta_e\mapsto g_i\Theta_eg_j^{-1}$, and only conjugacy classes of loop products are invariant
(\MT{prop:geo-edge-twist-invariance}). Nothing carries over to a continuum connection: graph-link
holonomy is neither base-connection holonomy nor bundle topology, and the technical text keeps these and
gauge redundancy apart.
\end{analogy}

\subsection{An inferred link field and its physical reading}

\Cref{comp:lie-links} changes the status of the links within this architecture. They are now latent
Lie-algebra generators under a declared normalized prior, rather than necessarily held background
values. Their Gaussian recognition stationarity balances prior stress against a matter current
(\MT{prop:link-field-equation}); for a comparison record that current has the outer-product form
of a delta rule (\MT{cor:link-delta-rule}). The generator history prior couples neighboring
episodes, and its linearization near zero has the form of a massive Euclidean vector operator, with
matter contributing a positive screening mass (\MT{prop:link-linearized-equation}). The central
equations and their approximation domain are explained in \cref{comp:lie-links} and their Gaussian
values in \cref{comp:gaussian-latent-links}.

\begin{analogy}{structural}
\textbf{A Proca-type operator on the interaction graph.} The cycle operator supplies a spatial
curl-type stiffness on graph copies, the history increment supplies a second difference in the
episode index, and the generator precision and matter response supply mass terms. This is the
form of a massive Euclidean field equation, not a derivation of a spacetime Proca field. The
episode index is an inference index. The most-probable-history inertia belongs to smoothing a
joint history; the Fisher recognition flow is first order in its own flow time
(\MT{prop:link-map-inertia}, \MT{prop:link-recognition-flow}).
\end{analogy}

The scaffolding verdict is unchanged: the base, principal bundle, frames, and reference point are
structural inputs that this inference does not construct (\MT{sec:phil-base}). Inferring links
within a declared model is a mathematical construction; recovering links from a target system's
records requires the separate operational evidence of \MT{open:phil-testability}. The held-link
model is recovered in the zero-width free-energy limit under its stated hypotheses
(\MT{thm:link-held-limit}), which connects the architectures without establishing a physical
action principle, an empirical field, or the formation mechanism below.

\subsection{Participation and why it does not yet distinguish the program}

Participation, in the technical sense, is reciprocal inferential dependence through records inside one
joint law, with no live posterior feeding a same-step generative factor (\MT{req:gen-typing-prohibition}).
It occurs in the same-time inference loop (block coordinate ascent, as in any connected hierarchical
model), in parameter learning (whose optimal prior over an unrestricted family is the aggregate
recognition law, \MTeq{eq:obs-aggregate-prior}), and in delayed feedback through a declared kernel (still
the free energy of one joint, \MT{prop:full-graph-episode-lifted-joint}). The verdict is blunt: none of
these separates the program from a hierarchical model with learned priors, and the resemblance to
Wheeler's participatory proposal \citep{Wheeler1990} is neither derived from it nor evidence for it.
Formation could separate them: collectives constituted by the coupled inference, held by a restoring
interaction, and dependent on the preceding state. Hence the formation theorems in this chapter.

\subsection{The four-cycle: a ring of agents choosing whom to listen to}

Four agents sit on a ring. Each has a scalar Gaussian belief (mean $\bar k_i$, unit variance) and a
scalar Gaussian model (mean $\bar m_i$, variance $1/2$); each edge carries a translation link $\omega_i$
in the abelian group $(\R,+)$. The cost of an edge is the transported relative entropy of its residuals.
Each agent listens to exactly one neighbor, and an edge carries energy only when an endpoint listens to
it (source gating). Of the sixteen listening assignments, the two perfect matchings, which split the ring
into two mutually attending pairs, are the candidate collectives. The kinetics is Langevin diffusion at
temperature $\mathrm T$ on the energy sheet of the current assignment, interrupted by a rate-one clock at
which every agent redraws whom it listens to from its posterior. Everything is symmetric under the
rotations and reflections of the ring.

\emph{Seeded formation.} From a Gaussian configuration and a uniformly random assignment, with small
probability the assignment is already a matching and the configuration lies near a seed point where every
edge is equally mismatched; the attended residuals then relax, the unattended ones grow, and two pairs
form. The probability is that of a lucky initial condition, independent of the small parameter and
unchanged if the prior learns (\MT{prop:phil-seeded-contrast-formation}).

\emph{Selection by slow learning.} Now let the listening prior learn slowly from the population's
recognition law. Its rest point is the recognition marginal, so learning has the form of the
expectation-maximization update of mixture weights, run in continuous time; averaged over the fast
kinetics, the listening probabilities climb the logarithm of an evidence function in the Fisher metric of
the simplex. At low temperature that evidence is largest at the two pure matchings and has a saddle at
the uniform prior, unstable exactly in the direction that favors one matching over the other. An
asymmetric Gaussian initial law leaves, during the early transient, an asymmetric imprint on the rows of
the order of the mobility; the saddle amplifies it, and after a time of order
$\epsilon_{\mathrm L}^{-1}\log(1/\epsilon_{\mathrm L})$ the prior sits at one matching while the
recognition law tracks the corresponding equilibrium. The frozen arm has nothing to amplify and relaxes
to its symmetric equilibrium (\MT{conj:phil-contrast-formation}).

\begin{analogy}{structural}
The theorem has the dynamical signature of mean-field symmetry breaking. The learned four-cycle is a
nonlinear Markov process, whose generator depends on its own law through the rows, as in the
McKean--Vlasov limit of the Curie--Weiss model; the rest point of learning is a Curie--Weiss-type
self-consistency (the listening marginal computed under the rows equals the rows), and below the
theorem's temperature threshold its symmetric solution is unstable along one direction. The limits fail
to commute in the familiar way: at fixed time the rows move by at most $\epsilon_{\mathrm L}t/2$ (part (a)
of \MT{prop:phil-seeded-contrast-formation}), so $\epsilon_{\mathrm L}\to0$ first and then $t\to\infty$
gives the symmetric frozen answer, while $t\propto\epsilon_{\mathrm L}^{-1}\log(1/\epsilon_{\mathrm L})$
gives a selected matching. The imprint, of the order of the mobility, acts as a vanishing field whose
effect survives because its amplification time diverges.

It is not spontaneous symmetry breaking in the strict sense. There is no thermodynamic limit, and the
frozen kinetics has a unique invariant law (the proof uses the Harris theorem for exactly this). Nothing
is spontaneous either: the rows learn from the law, so they are deterministic and carry no fluctuations
that could choose a branch, and the branch is fixed by a definite datum, the sign of an analytic function
of the initial mean. A symmetric initial law keeps the rows uniform forever (part (c) of
\MT{prop:phil-seeded-contrast-formation}), as a mean-field magnetization started at zero stays at zero.
The accurate description is explicit symmetry breaking by the initial data, amplified by a dynamical
instability; the technical text's word is ``amplifies''.
\end{analogy}

\begin{analogy}{structural}
Why matchings win at low temperature: every listening sheet has the same minimum energy, zero, at the
same configuration, so the sixteen sheets are degenerate. They differ in how many directions they
stiffen: a matching stiffens four residual directions, a path six, an oriented cycle seven (the residuals
around a loop obey one linear relation). Each extra stiff direction costs a factor of order
$\sqrt{\mathrm T}$ in thermal phase-space volume, so the evidence peaks at the matchings. Selection among
degenerate ground states by the volume of thermal fluctuations is the mechanism of order by disorder; the
technical documents do not use that term.
\end{analogy}

\emph{Late persistence, early null, and the parent.} Once the prior sits at a matching, the selected pairs
stay coherent, but the stationary law at those rows realizes the same persistence with probability
tending to one, and the frozen arm shows persistent pairs too, split half and half between the two
pairings (\MT{prop:phil-selected-pair-persistence}). Late, learning changes which pairing, not whether
there is pairing. Early, a slow prior has barely moved, and the law of the first formed block is the same
with or without learning (\MT{prop:phil-early-block-null}). The parent of a pairing is obtained by a fixed
channel that keeps the assignment, the loop coordinate $\sum_i\omega_i$, the slow directions, and one
fast mode (the sum of the two pairs' average residuals), and discards three fast modes (the difference of
the pairs' average residuals and each pair's belief-model residual difference). The parent is exact as a
summary when the discarded modes, given the kept ones, are thermal, and exact as a dynamics when its own
evolution reproduces the kept trajectories. Both hold after the late selection and, on the early seeded
episode, at deterministic activation times after the discarded modes have relaxed to the thermal scale;
the first moment the coherence criterion is met is too early by at least a factor of two in logarithmic
time, and there the discarded information diverges (\MT{prop:phil-early-first-completion}).

\begin{analogy}{exact}
The discarded conditional relative entropy is the error term of integrating out fast variables. The
chain rule $\KL(q\Vert\mu)=\KL(C_\#q\Vert C_\#\mu)+\mathcal H_{\mathrm{opt}}(q)$
(\MTeq{eq:phil-conditioned-parent-chain}) is an identity, so the fine divergence is the parent's
divergence plus a remainder that vanishes exactly when the fast modes are in conditional equilibrium given
the slow ones: adiabatic elimination written as an identity. The dynamical statements are the structural
analogue of a valid reduced Langevin equation, proved only on stated windows. Those windows lie on a
plateau between the thermalization time of the discarded modes, $\log(1/\epsilon)=\tfrac18\log(1/\mathrm
T)$, and the slow time $\mathrm T^{-1}$ on which weakly restored directions equilibrate and the pairing
can change; a quasi-stationary state between fast relaxation and slow escape is what a physicist calls
metastability.
\end{analogy}

\section{Objects and notation}
\label{comp:ch14-notation}

\begin{decoder}
Idle posit & A posit changing the law of no record or record statistic (\MT{def:phil-idle-wheel}). & Gauge redundancy; unobservable absolute frame. \\
$[\mathrm{Hol}^x(p)]_{\rm conj}$ & Conjugacy class of the product of links around closed walk $p$ in channel $x$. & Wilson loop before the trace. \\
Frame, law, mean agreement & Three inequivalent notions of agents agreeing (\MT{prop:geo-three-agreements}). & Pure gauge; equal states up to transport; equal first moments. \\
$\mathbf x=(\bar k,\bar m,\omega)$ & Four-cycle configuration: belief means, model means, link translations; twelve coordinates. & Site fields and abelian link fields. \\
$r^b_i$, $r^m_i$, $\mathcal D_{e_i}$ & Belief and model residuals on edge $i$ and the edge cost, a transported relative entropy. & Covariant difference and its energy. \\
$S$, $\sigma_i$ & Listening assignment; $\sigma_i=+1$ when $i$ listens to $i+1$. & Discrete label field. \\
$S^{(0)}$, $S^{(1)}$, matching & The two assignments that split the ring into two mutually attending pairs. & The two dimer coverings of a plaquette. \\
$\pi_{ij}$, $p_i=\pi_{i,i+1}$ & Learned source prior (listening probabilities), one row per receiver. & Mean-field order parameter. \\
$\epsilon_{\mathrm L}$, $\mathrm T$ & Learning mobility; temperature ($\mathrm T=\epsilon^8$ in the original regime). & Slow relaxation rate; temperature. \\
$q^{\epsilon_{\mathrm L}}_t$, $q^0_t$ & Recognition law of the kinetics, learned arm and frozen arm (rows fixed at $1/2$). & Ensemble with and without feedback. \\
$q^\star$, $\mu_{\mathrm T,p}$ & Restricted optimizer $\propto\pi_Se^{-U_S/\mathrm T}$, the invariant law at fixed rows $p$. & Gibbs state at fixed couplings. \\
$\nu^M_{\mathrm T}$ & Gaussian law on the sheet of matching $M$, assignment held at $M$. & Pure phase. \\
Contrast formation & A block internally coherent and externally distinct, entered after time zero. & Nucleation of a bound pair. \\
$\mathsf j_{\mathrm T}(\mathsf m_0)$ & Transient imprint: integrated early asymmetry of the frozen law along the matching direction. & Integrated response to an initial perturbation. \\
Restored equilibrium coherence & Coherence with a restoring interaction that is typical at stationarity; not formation. & Equilibrium order. \\
$C$, $W_M$, $Z_M$ & Parent channel; kept subspace (nine dimensions); discarded subspace (three fast modes). & Blocking map; kept and integrated-out variables. \\
$\mathcal H_{\mathrm{opt}}$, $\mathcal H_{\mathrm{phys}}$ & Mean discarded conditional relative entropy against the optimizer or the physical posterior. & Error of adiabatic elimination. \\
Retained-path TV & Total variation between the parent's path law and the fine kept-coordinate path law. & Validity of a reduced dynamics. \\
$\theta$, $\mathsf t_{\mathrm c}$ & Deterministic activation time; first qualifying completion time. & Fixed measurement time; first-passage time. \\
\end{decoder}

Several words carry nonstandard senses. ``Episode'' means a coherence episode in the sense of rule R15,
not a learning episode. ``Formation'' is always a threshold event at declared resolutions (relative
entropies below $\epsilon^4$ inside, above $1/2$ across the boundary, with an entrance clause), and
``selection'' always a property of the law, never of a sample path. ``Actual'' means the law the kinetics
really produces, not a Gaussian or stationary substitute. ``Physical'' refers to the generative posterior,
``optimizing'' to the restricted optimizer. The regime parameter appears as $\epsilon$ in the appendices,
$\xi$ in the chapter text, and $\mathrm T^{1/8}$ in \MT{prop:phil-early-block-null}; the original coupled
regime sets the mobility $\epsilon_{\mathrm L}$ equal to it. The label
\texttt{conj:phil-contrast-formation} survives from when the statement was a conjecture; the printed
heading says Theorem, and it is proved.

\section{The formulas, explained}
\label{comp:ch14-formulas}

\subsection{Energy, invariant law, and learning rule}

On edge $e_i$ the residuals are $r^b_i=\bar k_{i+1}-\bar k_i-\omega_i$ and
$r^m_i=\bar m_{i+1}-\bar m_i-\omega_i$, and
\[
\mathcal D_{e_i}(\mathbf x)=\tfrac12\bigl[(r^b_i)^2+(r^m_i)^2\bigr]
=\KL\bigl(q^b_{i+1}\Vert(\Theta_{e_i})_\#q^b_i\bigr)
+\sigma^2_{\mathrm{ev}}\KL\bigl(q^m_{i+1}\Vert(\Theta_{e_i})_\#q^m_i\bigr)
\]
(\MTeq{eq:rg-formation-cycle-cost}): the relative entropy between one agent's law and its neighbor's law
carried across the link. The sheet energy $V_S=\sum_i\mathcal D_{e(i,S_i)}$ counts an edge twice if both
ends attend to it and not at all if neither does. The restricted optimizer, which is the invariant law of
the kinetics at a fixed prior, is
\[
q^\star(\mathbf x,S)\propto\pi_Se^{-U_S(\mathbf x)/\mathrm T},
\qquad
U_S=V_S+\tfrac{\mathrm T}2|\mathbf x|^2+\tfrac{\mathrm T}2\mathbf x^\top R_{\mathrm{ev}}\mathbf x.
\]
The first term is the interaction; the second comes from the unit Gaussian prior on configurations and
gives every direction a weak restoring rate of order $\mathrm T$; the third comes from the model evaluator
and penalizes belief-model discrepancy, and the physical posterior omits it. At zero temperature every
sheet has a flat zero-energy valley; at finite temperature each sheet is a Gaussian whose weight is its
phase-space volume.

The learning rule of the source prior is
\[
\dot\pi_{ij}=\epsilon_{\mathrm L}\bigl(q_t(S_i=j)-\pi_{ij}\bigr)
\]
(\MTeq{eq:phil-source-learning}): the rows relax at rate $\epsilon_{\mathrm L}$ toward the current
recognition marginal. It is the natural-gradient flow of the free energy in the prior row under the Fisher
metric of the simplex, with the recognition label marginal as rest point, as for the aggregate prior. The
right side uses the law $q_t$, not a sample, so the rows are deterministic; at $\epsilon_{\mathrm L}=0$
they stay at $1/2$, the frozen arm against which every selection claim is measured. Part (a) of
\MT{prop:phil-seeded-contrast-formation} gives well-posedness and $|\pi_{ij}(t)-1/2|\le\epsilon_{\mathrm
L}t/2$, so nothing learned can matter on horizons short compared with $1/\epsilon_{\mathrm L}$.

\subsection{Where the probability of formation comes from}

Part (b) of \MT{prop:phil-seeded-contrast-formation} bounds the probability of contrast formation below by
\[
\frac{e^{-1/4}}{16}\sum_{\pm}\mathcal N(\mathsf m_0,I_{12})
\bigl\{|\mathbf x\mp\mathbf x^\circ|<\varrho_\circ\bigr\},
\qquad \mathbf x^\circ=(0,0,\mathbf 1_4),\quad \varrho_\circ=\tfrac1{100},
\]
uniformly in the regime parameter. The factor $1/16$ is the probability that the uniform initial
assignment is a given matching; $e^{-1/4}$ is the probability that the redraw clock is silent until time
$1/4$, long enough for a cost gap to open; the Gaussian mass is the probability of starting within $1/100$
of a seed point, where all means vanish and every link translation is $\pm1$, so every edge is equally
mismatched. On this sub-event the mechanism of \MT{prop:rg-formation-endogenous-cycle} runs: attended
residuals decay at rates of order one, unattended ones grow toward $3/2$ in magnitude, and posterior
redraws keep the matching except with exponentially small probability. The bound holds equally with
frozen rows, and the technical text draws the right conclusion: a claim that learning creates a contrast
must compare the arms at the level of the law. Two calibrations: the explicit estimate
$e^{-1/4}\varrho_\circ^{12}/737280$ times a Gaussian factor is about $10^{-31}$ at $\mathsf m_0=0$, so its
only role is to be positive and $\epsilon$-independent; and the seeded sub-event is genuine formation
under R15, with a restoring interaction and a probability exceeding its value from the invariant law, at
most $4e^{-1/(4\mathrm T)}$, by a margin bounded below.

\subsection{The row flow, its saddle, and the imprint}

Averaging the learning rule over the invariant law at fixed rows gives the reduced field
\[
\mathscr V_i(p)=\bar\beta_i(p)-p_i=p_i(1-p_i)\partial_{p_i}\log\mathsf Z(p),
\qquad
\mathsf Z(p)=\sum_\sigma\pi_\sigma(p)\mathsf w_\sigma,
\qquad
\mathsf w_\sigma=\det(\mathsf A_\sigma/\mathrm T)^{-1/2}
\]
(\MTeq{eq:phil-selection-reduced-field}). Here $\bar\beta_i(p)$ is the stationary probability that
receiver $i$ listens to $i+1$, $\mathsf w_\sigma$ is the Gaussian volume of sheet $\sigma$, and
$\mathsf Z$ is the evidence of the prior. The factor $p_i(1-p_i)$ is the inverse Fisher metric of a
Bernoulli row, so the reduced flow is the replicator equation with fitness $\partial_{p_i}\log\mathsf Z$,
natural-gradient ascent of the log evidence. With the matching weight normalized to one, paths have
relative weights $\mathsf r_E,\mathsf r_N\approx\sqrt{40/7}\mathrm T$ and oriented cycles
$\mathsf r_C\approx\sqrt{5/3}\mathrm T^{3/2}$, the count of stiff directions made exact. At the uniform
prior the linearized field has one positive eigenvalue $\alpha\ge7/5$ along $\mathbf v=(1,-1,1,-1)$,
toward one matching and away from the other, and negative eigenvalues elsewhere
(\MTeq{eq:phil-selection-saddle-spectrum}); both matching vertices are sinks. The threshold
$\mathrm T_0$ is a temperature below which the sheet ratios are at most $1/8$: a sufficient condition of
the proof, not a located critical temperature, and nothing is proved above it.

The seed that pushes the rows off the saddle is the transient imprint
\[
\mathsf j_{\mathrm T}(\mathsf m_0)=\mathbf v\cdot\mathsf I(\mathsf m_0),
\qquad
\mathsf I_i(\mathsf m_0)=\int_0^\infty\bigl[q_t^0(\sigma_i=+1)-\tfrac12\bigr]dt
\]
(\MTeq{eq:phil-selection-imprint}), the time-integrated early asymmetry of the frozen law. It is an entire
function of the initial mean and changes sign under rotation of the ring by one site, so it vanishes at
every rotation-invariant mean, the seed points included. The proof shows it is not identically zero, so
its zero set is Lebesgue-null: that is the excluded set of means. After the transient the rows sit at
$1/2+\epsilon_{\mathrm L}\mathsf I(\mathsf m_0)+o(\epsilon_{\mathrm L})$; the component along
$\mathbf v$ grows like $\epsilon_{\mathrm L}e^{\alpha\epsilon_{\mathrm L}t}$ and reaches order one at
$t\approx\alpha^{-1}\epsilon_{\mathrm L}^{-1}\log(1/\epsilon_{\mathrm L})$, hence the window
$t_{\mathrm L}=\mathrm C_{\mathrm T}\epsilon_{\mathrm L}^{-1}\log(1/\epsilon_{\mathrm L})$ with
$\mathrm C_{\mathrm T}>1/\alpha$. The theorem states: there is $\mathrm T_0>0$ such that for every
$\mathrm T\in(0,\mathrm T_0]$ and Lebesgue-almost every $\mathsf m_0$ one matching $M$ satisfies
\[
\lim_{\epsilon_{\mathrm L}\to0}\ \sup_{t_{\mathrm L}\le t\le2t_{\mathrm L}}
q^{\epsilon_{\mathrm L}}_t(S\ne M)=0,
\]
while for every $\mathrm T>0$ and every $\mathsf m_0$ the frozen difference
$q^0_t(S=S^{(0)})-q^0_t(S=S^{(1)})$ tends to zero as $t\to\infty$. The order of limits is fixed
temperature, then mobility to zero, with the window scaled by the mobility; the conclusion concerns the
one-time law of the assignment on a moving window, not histories, events, or $t\to\infty$ at fixed
mobility.

\subsection{What the corollary and the tracking bound add}

\MT{cor:phil-joint-selection} lets the temperature vanish too. Along any $\mathrm T_n\downarrow0$ with
tolerances $\eta_n\to0$ there are mobility thresholds $d_n\downarrow0$, the same for almost every mean,
below which the selection error on the window (now with $\mathrm C_{\mathrm T}=2$ for all temperatures)
is at most $\eta_n$, the learned law is $\eta_n$-close in relative entropy to its instantaneous
equilibrium, and the frozen law is $\eta_n$-close to its symmetric equilibrium, so the frozen matching
probabilities tend to $1/2$. It adds a single joint path $(\mathrm T(\epsilon_{\mathrm L}),
\epsilon_{\mathrm L})\to0$ and an assignment-law contrast tending to $1/2$ in total variation. It adds no
rate: the thresholds exist by a measure-theoretic argument and have no values, they can be taken below
$\mathrm T_n^a$ for any $a>7/2$ only as a cap, and the matching may change with $n$, because nothing
controls the imprint's sign as the temperature vanishes.

\MT{prop:phil-temperature-tracking} makes the adiabatic half explicit. For $\mathrm T\le1/100$ the frozen
entropy production dominates $\mathrm c_*\mathrm T^{5/2}/[4+3\log(1/\mathrm T)+8\log(1/p_{\min})]$
times the relative entropy, and
\[
\operatorname{TV}(q_t,\mu_{p(t)})
\le e^{-\ell t/2}\sqrt{\frac{13(12+|\mathsf m_0|^2)}{4\mathrm T}}
+\frac{2\epsilon_{\mathrm L}\mathrm D(\mathrm T,\epsilon_{\mathrm L})}{\mathrm c_*\mathrm T^{7/2}},
\qquad
\ell=\frac{\mathrm c_*\mathrm T^{5/2}}{\mathrm D(\mathrm T,\epsilon_{\mathrm L})}
\]
(\MTeq{eq:phil-temperature-tracking}), with $\mathrm D$ logarithmic. The first term is memory of the
initial condition, erased at the mixing rate $\ell$; the second is the lag of a law chasing an equilibrium
that moves at a rate of order $\epsilon_{\mathrm L}/\mathrm T$, divided by the mixing rate, which is where
the exponent $7/2$ comes from. For $\epsilon_{\mathrm L}=\mathrm T^a$, $a>7/2$, tracking holds on the
window: an explicit adiabatic condition, not a selection result, since whether the equilibrium has moved
to a matching depends on the imprint, for which no lower bound as $\mathrm T\to0$ is available.

\subsection{Persistence, and why it is equilibrium}

\MT{prop:phil-selected-pair-persistence} assumes the learned law is $\delta$-close to $\nu^M_{\mathrm T}$
and the frozen law $\delta$-close to its symmetric equilibrium. Over a window $\mathsf h$ the learned path
law is then within $(\mathsf h+2)\delta$ of a stationary Ornstein--Uhlenbeck path held at $M$, the frozen
one within $\delta+16\mathrm T+31\mathrm T\mathsf h$ of an equal mixture of the two held-matching paths,
and coherence at resolution $\mathsf a$ fails with probability at most
\[
(\mathsf h+2)\delta+\mathsf b(\mathrm T,\mathsf a,\mathsf h),
\qquad
\mathsf b=8\sqrt{2/\pi}\mathsf a+(8+3400\mathsf h)\mathrm T/\mathsf a^2
\]
(\MTeq{eq:phil-selected-coherence-persistence}). The first term of $\mathsf b$ is the chance that an
exterior residual, of variance of order one, is accidentally below $\mathsf a$; the second is the chance
that thermal noise pushes an internal residual above $\mathsf a$ during the window. The choice
$\mathsf a=\mathrm T^{1/4}$, $\mathsf h=\mathrm T^{-1/4}$ puts the resolution far above the thermal scale
(restoring ratio $4\mathsf a^2/\mathrm T\to\infty$) and the window far below the hopping time
$\mathrm T^{-1}$; the error is then of order $3400\mathrm T^{1/4}$, small only far below
$\mathrm T=10^{-14}$. The last clause decides the classification: with rows near the matching vertex, the
stationary law at those frozen rows realizes the same persistence event with probability at least
$1-4(\mathsf h+2)\delta-\mathsf b$. The event is typical at equilibrium, fails the horizon clause of R15,
and is restored equilibrium coherence. The coherent-partition laws of the two arms differ by nearly $1/2$
because the frozen arm splits its equilibrium pairs evenly (\MTeq{eq:phil-coherent-partition-contrast}).

\subsection{The early first-block null}

\MT{prop:phil-early-block-null} keeps the original contrast detector on the horizon
$H_{\mathrm T}=2\mathrm T^{-1/8}$ and takes the block that completes first. Coupled with common noise,
clock, and initial sample, the arms differ only when a redraw differs, and a row displaced by
$\epsilon_{\mathrm L}t/2$ changes a redraw probability by at most $\epsilon_{\mathrm L}t$. The conditional
first-block laws differ by at most $\mathsf d_{\mathrm T}/\mathsf c(\mathsf m_0)$, with
$\mathsf d_{\mathrm T}=\min\{1,8\epsilon_{\mathrm L}\mathrm T^{-1/4}\}$ and $\mathsf c$ the seeded bound
above (at most of order $10^{-30}$), so they agree when $\epsilon_{\mathrm L}=o(\mathrm T^{1/4})$, at
extraordinarily small temperatures; at the original mobility $\mathrm T^{1/8}$ the bound is vacuous.

\subsection{Conditioning, and what a vanishing discarded term means}

The late parent starts from the full recognition law conditioned on a completed coherence prefix $E$ on
$[t_{\mathrm L},t_{\mathrm L}+1]$, pushed through the channel. \MT{prop:phil-actual-conditioned-parent}
bounds the conditioned relative entropy and the optimizing discarded term by
\[
\kappa_E=\frac{\eta+h_b(\rho)}{1-\rho},
\qquad
h_b(r)=-r\log r-(1-r)\log(1-r),
\]
and the physical discarded term by $\kappa_E+\mathrm T$, where $\eta$ bounds the unconditioned relative
entropy to equilibrium and $\rho$ the failure probability of the prefix. The binary entropy is the price
of conditioning, small only for a typical event. \MT{rem:phil-conditioned-hidden-sign} shows the price is
real: conditioning a thermal Gaussian on the sign of one discarded coordinate costs $\log2$ at every
temperature, and conditioning on the whole realized history pins the discarded coordinates and makes the
relative entropy infinite. A parent is activated from a law given a high-probability event, never from
one history. With $\eta,\delta\le\mathrm T$ and $\epsilon_{\mathrm L}\le\mathrm T^2$,
$\kappa_E=O(\mathrm T^{1/4}\log(1/\mathrm T))$.

\MT{prop:phil-actual-parent-dynamics} compares paths. The parent diffuses on the kept subspace and, at each
ring of its own clock, reinserts the discarded coordinates as fresh thermal draws, redraws the assignment
from the posterior at the reconstructed configuration, and projects back; its rows learn from its own law.
Its path law, including the assignment history and hence the loop, is within
\[
2\rho+(\mathsf h+3)\delta+\frac{4\mathsf hd}{1-4d},
\qquad d=\delta+\epsilon_{\mathrm L}\mathsf h,
\]
of the conditioned fine kept-coordinate path law (\MTeq{eq:phil-actual-parent-path-error}), which is
$O(\mathrm T^{1/4})$ for $\mathsf h=\mathrm T^{-1/4}$: the cost of conditioning, the distance of both
processes from the held-matching stationary path, and the probability that the parent leaves the matching.

\subsection{The activation clock}

On the early seeded episode, at $\mathrm T=\epsilon^8$ and any $\epsilon_{\mathrm L}\in[0,1]$,
\MT{prop:phil-early-seeded-parent} bounds the discarded information at a deterministic activation time
$\theta$ from both sides. Asymptotically,
\[
\mathcal H_{\mathrm{opt}}(q^\circ_\theta)\le\tfrac92\varrho_\circ^2e^{-8(\theta-\log(1/\epsilon))}+o(1),
\qquad
\mathcal H_{\mathrm{opt}}(q^\circ_\theta)\gtrsim\frac{\varrho_\circ^2}{7}\mathrm c_{\mathsf m}
e^{8(\log(1/\epsilon)-\theta)}-\log2
\quad(\theta\le\log(1/\epsilon)),
\]
with $q^\circ_\theta$ the law conditioned on the seeded prefix and $\mathrm c_{\mathsf m}>0$ depending
only on the initial mean. The physics is one line. The discarded discrepancy coordinates start at offsets
of order $\varrho_\circ$ and decay at rate about $4$, while their thermal width is of order
$\sqrt{\mathrm T}=\epsilon^4$; a displaced Gaussian has relative entropy of order squared offset over
variance, $\varrho_\circ^2e^{-8\theta}/\mathrm T=\varrho_\circ^2e^{-8(\theta-\log(1/\epsilon))}$. So the
discarded information is large before the thermalization time $\log(1/\epsilon)$ and small after it. The
upper bound needs $\theta-\log(1/\epsilon)\to\infty$ and $\epsilon^4\theta\to0$, the lower bound diverges
when $\log(1/\epsilon)-\theta\to\infty$, and the window $\theta=\log(1/\epsilon)+O(1)$ is undecided. The
original witness of \MT{prop:rg-formation-endogenous-cycle} activates at $2\log(1/\epsilon)$, safely on
the good side.

\MT{prop:phil-early-first-completion} treats the clock a physicist might pick first: the first time any
block meets the contrast criterion. That criterion asks for residuals below about $\epsilon^2$, far
coarser than the thermal scale $\epsilon^4$. Seeded residuals fall below $\epsilon^2$ by about
$\tfrac12\log(1/\epsilon)$, and near $\tfrac14\log(1/\epsilon)$ when a pair's initial discrepancy is small,
so the first completion occurs by $\theta_\circ=\tfrac12\log(1/\epsilon)+1$. There the discarded
information is at least of order $\mathrm T^{-5/8}$, driven by seeds whose first pair completes near
$\tfrac14\log(1/\epsilon)$ while the other matched pair, whose discrepancy the channel also discards,
still carries an imprint of order $\varrho_\circ\epsilon$. The divergence persists when the law is
conditioned on the whole contrast event. This is an established negative result: no parent of this type
activated at the first qualifying completion has a vanishing conditional error, in either arm and at any
mobility in $[0,1]$.

\MT{prop:phil-early-parent-dynamics} bounds the path error of a parent activated at a good deterministic
time by
\[
24\exp\Bigl[-\frac1{76800\mathrm T\mathsf h}\Bigr]
+4\bigl(1+\sqrt{8/3}\bigr)\Lambda_{\mathrm L}(\theta+\mathsf h)\mathsf he^{-1/\mathrm T},
\qquad
\Lambda_{\mathrm L}(t)=2e^{\epsilon_{\mathrm L}t}-1,
\]
uniformly in the mean and in $\epsilon_{\mathrm L}\in[0,1]$, with $\Lambda_{\mathrm L}$ bounding how far
learned rows can tilt a redraw. The first term is the probability that thermal diffusion over the window
moves the configuration a distance $1/20$ off its deterministic path, negligible when
$\mathrm T\mathsf h\to0$; the second is the Arrhenius probability of a redraw across a cost gap larger than
one. Windows up to $\mathrm T^{-1+\mathrm q}$, $\mathrm q\in(0,1)$, are covered, and nothing is concluded
at order $\mathrm T^{-1}$. The late proof cannot be reused here because the early activation law is at
total-variation distance essentially one from the stationary kept-coordinate Gaussian
(\MT{rem:phil-early-stationary-mismatch}), so the early proof follows the energy gap along the
conditioned path.

\section{What is proved, and how solid it is}
\label{comp:ch14-scorecard}

The first scorecard covers the interpretive layer.

\begin{scorecard}
\MT{def:phil-idle-wheel} & A posit that changes the law of no record or record statistic is removed. & \CHOICE & Declared rule of reading, not derived from constructive empiricism. \\
Local-first verdict, \MTeq{eq:phil-fiber-measures} & Agents' law points are primary; base, bundle, spaces are scaffolding with no law. & \CHOICE & Describes the present architecture; changes if edges become inferred coordinates. \\
\MTeq{eq:phil-invariants} & Evidence values, loop classes, and per-edge twists up to conjugation survive frame change. & \SOLID & Invariance proved in the cited statements; reading it as ``structure'' is interpretation. \\
Realist and Neo-Kantian readings & Compatible with the invariants; none selected. & \VAGUE & The chapter declines to select and names what support would need. \\
\MT{prop:link-field-equation}, \MT{prop:link-linearized-equation} & Latent links have a stationarity equation and a Proca-type linearized operator. & \SOLID & Declared prior, stated matter and history hypotheses; physical-field reading remains structural. \\
\MT{thm:link-held-limit} & Held-link free-energy infima and evidence are zero-width limits of latent-link quantities. & \COND & Bounded records, positivity, and the theorem's continuity and variational hypotheses. \\
\MT{hyp:phil-observational-closure} & All data reach an agent as clamped states or node messages. & \CHOICE & Free for kernels by \MT{thm:obs-agent-interaction-equivalence}; removes no conditioning. \\
\MT{hyp:phil-boundary-agents} & Boundary environment nodes are called agents. & \CHOICE & Naming only; adds no state, action, blanket, or free energy. \\
Participatory reading & Participation via inference loop, learning, delayed feedback; formation could distinguish the program. & \VAGUE & Negative half rests on proved identifications; positive half is programmatic. \\
\MT{claim:physical-law-identification} & Identify a coarse-graining object with a physical law. & \OPEN & Needs a named target system and held-out evidence; none supplied. \\
\MT{open:phil-testability} & An observation at $c_*$ that could count against the program. & \OPEN & Only model-internal separations are proved. \\
\end{scorecard}

The second covers the continuation line. Every theorem-type statement is proved only on the four-cycle,
with abelian translation links on scalar Gaussian means, so all are TOY; the two general items are marked
separately.

\begin{scorecard}
\MT{prop:phil-seeded-contrast-formation} & Contrast formation from every Gaussian initial law with probability bounded below; same with frozen rows; symmetric laws keep rows uniform. & \TOY & Probability from the matching seed; learning idle; floor about $10^{-31}$ at $\mathsf m_0=0$. \\
\MT{conj:phil-contrast-formation} & Fixed small $\mathrm T$, $\epsilon_{\mathrm L}\to0$: learned assignment law concentrates on one matching on $[t_{\mathrm L},2t_{\mathrm L}]$ for a.e.\ mean; frozen arm symmetric. & \TOY & Amplifies the initial-law asymmetry; one-time law; no rate in $\mathrm T$. \\
\MT{cor:phil-joint-selection} & Joint limit $\mathrm T,\epsilon_{\mathrm L}\to0$ along implicit thresholds; frozen matching probabilities tend to $1/2$. & \TOY & Thresholds existential; matching may vary with $\mathrm T_n$. \\
\MT{prop:phil-temperature-tracking} & Entropy-production rate of order $\mathrm T^{5/2}$; tracking for $\epsilon_{\mathrm L}=\mathrm T^a$, $a>7/2$. & \TOY & Tracking only; no power-law selection. \\
\MT{prop:phil-selected-pair-persistence} & Selected pair coherent at resolution $\mathrm T^{1/4}$ for time $\mathrm T^{-1/4}$; partition contrast near $1/2$. & \TOY & Restored equilibrium coherence; no entrance claim. \\
\MT{prop:phil-early-block-null} & For $\epsilon_{\mathrm L}=o(\mathrm T^{1/4})$ learned and frozen first-block laws agree in the limit. & \TOY & Silent at the original mobility $\mathrm T^{1/8}$. \\
\MT{prop:phil-actual-conditioned-parent} & Late activation: relative entropy and discarded information $O(\mathrm T^{1/4}\log(1/\mathrm T))$. & \TOY & Averaged over a declared prefix, not per history. \\
\MT{prop:phil-actual-parent-dynamics} & Parent path law within $O(\mathrm T^{1/4})$ of the conditioned fine kept path over $\mathrm T^{-1/4}$. & \TOY & Finite window; own-law learning; no recursion. \\
\MT{prop:phil-early-seeded-parent} & Discarded information vanishes if $\theta-\log(1/\epsilon)\to\infty$, diverges if $\log(1/\epsilon)-\theta\to\infty$. & \TOY & Seeded sub-event; window $\log(1/\epsilon)+O(1)$ open. \\
\MT{prop:phil-early-first-completion} & At first qualifying completion, discarded information at least of order $\mathrm T^{-5/8}$, also given the whole contrast event. & \TOY & Established negative result on the four-cycle. \\
\MT{prop:phil-early-parent-dynamics} & Early parent path error vanishes when $\mathrm T\mathsf h\to0$, uniformly in mean and mobility. & \TOY & Windows shorter than $\mathrm T^{-1}$. \\
\MT{lem:phil-early-discarded-information} & Discarded information under conditioning, mixing, and extra variables; a Donsker--Varadhan lower bound. & \SOLID & General; standard in substance (chain rule). \\
\MT{rem:phil-conditioned-hidden-sign} & Conditioning on a probability-$1/2$ event costs $\log2$; on a realized history, infinity. & \SOLID & General counterexample shown on a Gaussian. \\
\end{scorecard}

\begin{hygiene}
What matters is what makes the four-cycle special and what fixes the comparison. The graph has exactly two
perfect matchings; links are abelian translations on scalar means, so loop holonomy is the number
$\sum_i\omega_i$ and plays no role in any selection; edges carry energy only when attended. The temperature
is small and fixed ($\mathrm T\le\mathrm T_0$ or $1/100$ late, $\mathrm T\le10^{-16}$ for the early null,
$\mathrm T=\epsilon^8$ with $\epsilon\le1/100$ for the early parent), and the bounds are asymptotic, some
with large constants. The rows learn from the law, not from samples, which makes them deterministic and
excludes fluctuation-driven selection: a modeling commitment, not hygiene. The initial law is Gaussian
with full support, and ``almost every mean'' excludes the imprint's zero set, which contains every
rotation-invariant mean. The kinetics are a declared realization along which the free energy does not
increase, not selected by it; the parent channel is fixed by hand. The late results need
$\eta,\delta\le\mathrm T$ and $\epsilon_{\mathrm L}\le\mathrm T^2$; the early ones condition on the seeded
sub-event, not the whole contrast event. Hygiene to take on trust: standard Borel spaces, measurability
of the first-block mark, well-posedness, finite moments, density envelopes, and form-domain extensions of
functional inequalities. Observational closure costs nothing for kernels; the boundary-agent hypothesis is
a naming decision.
\end{hygiene}

\section{Toy versus general}
\label{comp:ch14-toy}

\subsection{Mechanism and small-network specifics}

Every formation statement here lives on four agents; the program's target is 512 or more agents with
nonabelian or non-compact groups such as $\GL(K,\R)$ or $\mathrm{SU}(N)$. The four-cycle grounds and
validates; it does not set direction, and none of its results transfers without proof.
\Cref{comp:tab-ch14-toy} separates plausible mechanism-level content, the companion's reading and
conjectural beyond four agents, from four-agent specifics.

\begin{sciencetable}{Mechanism-level content versus four-cycle specifics.}
\label{comp:tab-ch14-toy}
\small
\begin{tabularx}{\textwidth}{>{\raggedright\arraybackslash}p{0.16\textwidth}>{\raggedright\arraybackslash}X>{\raggedright\arraybackslash}p{0.31\textwidth}}
\toprule
\textbf{Result} & \textbf{Mechanism-level content (conjectural beyond four agents)} & \textbf{Four-cycle specifics}\\
\midrule
Seeded formation & A favorable initial condition relaxes into a coherent block; the probability is the measure of such conditions. & Seed points $\pm\mathbf x^\circ$, floor near $10^{-31}$, factors $1/16$, $e^{-1/4}$. \\
Symmetric laws & A unique equivariant law-level dynamics cannot break a symmetry of its initial data; elementary and general, stated only for the four-cycle. & The form $\mathsf m_0=(\mathrm c_1\mathbf 1_4,\mathrm c_2\mathbf 1_4,0)$. \\
Pairing selection & Slow law-level learning is natural-gradient ascent of the log evidence; low temperature favors sheets with fewest stiff directions; a transient imprint decides the branch after time $\epsilon_{\mathrm L}^{-1}\log(1/\epsilon_{\mathrm L})$. & Two matchings, one unstable direction, ratios $\sqrt{40/7}\mathrm T$, $\sqrt{5/3}\mathrm T^{3/2}$, $\alpha\ge7/5$, the imprint computation. \\
Tracking & Mobility must be small against the frozen mixing rate. & Rate $\mathrm T^{5/2}$, exponent $7/2$. \\
Persistence & A structure selected in the prior is coherent in equilibrium on windows short against the hopping time. & Resolution $\mathrm T^{1/4}$, window $\mathrm T^{-1/4}$, flux $31\mathrm T$. \\
Early null & Learning slower than formation cannot bias the first block. & Threshold $o(\mathrm T^{1/4})$, horizon $2\mathrm T^{-1/8}$. \\
Parent closure & Exact to the extent discarded modes are in conditional equilibrium; activation must follow their thermalization; a detector coarser than the thermal scale fires too early. & Time $\log(1/\epsilon)$, divergence $\mathrm T^{-5/8}$, windows $\mathrm T^{-1/4}$, $\mathrm T^{-1+\mathrm q}$. \\
\bottomrule
\end{tabularx}
\end{sciencetable}

Some hypotheses exclude the target outright. Scalar fibers mean $K=1$, and abelian links mean holonomy
cannot act, so whether holonomy selects components is not touched; it is one branch of
\MT{open:rg-formation-adaptive-family}. The stationary preference for the fewest stiff directions does
generalize within the gated family: minimum edge covers dominate the stationary law on any finite graph
(\MT{prop:rg-gated-depth-one}). But nothing is proved about which of the combinatorially many covers slow
learning selects, or whether hopping times between covers grow with $N$, which is where genuine
metastability and a large-$N$ transition could enter; the four-cycle rules none of this in or out.
Recursion is out of reach too: eligible blocks have two or three agents, so the pairs cannot coalesce
within the event class, and recursive formation in the sense of R14 is posed outside the gated family
(\MT{open:rg-formation-general}).

\subsection{What the four-cycle says about waiting out the transient}

The reader's intuition is that meta-agent formation is soft and continuous: wait out the transient and
keep the slow-to-relax mode. The four-cycle gives this a precise form at $N=4$ in three respects and
qualifies it in two.

The activation-clock results are the sharpest support. A parent on the seeded episode is exact, as a
summary and as a dynamics, when activated at a deterministic time after the discarded modes have
thermalized, and provably not at the first moment a threshold criterion is met
(\MT{prop:phil-early-seeded-parent,prop:phil-early-first-completion}). ``Wait out the transient'' is
exactly $\theta-\log(1/\epsilon)\to\infty$, and the negative result says that a first-passage clock
driven by a detector coarser than the thermal scale is the wrong clock. The waiting is bounded on the
other side too: valid activation times and windows lie below the slow time $\mathrm T^{-1}$, so the
meta-agent is a quasi-stationary description on a plateau, not a permanent object. Second, what survives
the transient late is equilibrium structure: the selected pair is coherent because the equilibrium at the
learned prior is paired. Third, the channel does keep the slow directions, relaxing at rates of order
$\mathrm T$, with the loop and the assignment, and every discarded coordinate is fast.

The first qualification is that the channel was not chosen by a slow-mode criterion: it is supplied, it
keeps one fast mode as well, and no statement derives it from relaxation rates. The second cuts deeper.
Waiting out the transient selects nothing by itself. In a symmetric system with law-level learning, which
pairing forms is a property of sample paths, and at the level of the law nothing is selected unless a
datum breaks the symmetry (part (c) of \MT{prop:phil-seeded-contrast-formation}, rule R13). Soft
continuous formation and law-level selection of a culture are different claims; on the four-cycle the
first holds and the second needs an asymmetric initial law. Whether large networks add a genuine source
of selection, such as branch choice by fluctuations under learning from finite samples, or a transition
with diverging hopping times, nothing here addresses.

\section{Open ends}
\label{comp:ch14-open}

Physical-law identification (\MT{claim:physical-law-identification}) asks for a named target system whose
data the model does not generate, a record kernel separating two coarse-graining objects that differ in a
gauge-invariant datum, and held-out evidence, robust across blocking schemes, that a baseline margin is or
is not met. Single-point testability (\MT{open:phil-testability}) asks for a test at $c_*$, fixed before
the data, whose failure would falsify nontrivial holonomy on a walk or formation of a block, or a proof
that no record statistic distinguishes the alternatives. What exists is model-internal (the holonomy
record of \MT{thm:cg-loop-operational-holonomy}, the learned-versus-frozen selection); no target system,
noise model, estimator, or baseline is supplied.

The central formation open end is \MT{open:rg-formation-adaptive-family}: at the original coupled mobility
$\epsilon_{\mathrm L}=\mathrm T^{1/8}$, a difference bounded below between the learned and frozen
conditional laws of the first formed block, or between two initial means differing only in their loop
coordinate, with parents of vanishing errors; or a proof that these differences vanish for almost every
mean. The late selection and the early null bracket this regime without entering it. Also undecided are
the activation window $\theta=\log(1/\epsilon)+O(1)$, the discarded information given the whole contrast
event at a deterministic time, windows of order $\mathrm T^{-1}$ or longer, and a power-law selection
regime, which needs a quantitative lower bound on the imprint. Beyond the four-cycle, repeated endogenous
formation (\MT{open:rg-formation-general}), dynamic closure of the pointwise parent
(\MT{open:rg-pointwise-parent-dynamics}), edge creation (\MT{open:obs-edge-creation}), and a natural
partition selector (\MT{open:cg-natural-selector}) remain open; the last two mark where the local-first
verdict would change.

\begin{verdictbox}{\NARROWER}
The technical chapter says its readings are definitions, hypotheses, and open problems, and that it proves
a proposition and a theorem that sharpen the participatory reading: formation from every Gaussian initial
law with learning idle, and law-level pairing selection by slow source learning against a symmetric
frozen arm. Every proved statement does exactly what it says, the scope paragraph is accurate, and the
interpretive verdicts (local-first, no realist position selected, no agents created by naming) correctly
describe the formalism. For the claim that matters to the program, the result is narrower than it sounds.
What is selected is one of two pairings of a four-agent ring with abelian links, in the one-time
assignment law, by amplifying an asymmetry of the initial law, at fixed small temperature with mobility
tending to zero. The selected pair's later coherence is equilibrium coherence, not formation; the early
formation that passes the formation test is seeded and unaffected by slow learning; parent closure is
proved as averaged, finite-window statements on declared prefixes, and the first-completion clock is
refuted. The chapter itself concedes that assignment selection does not establish the recursive
participatory claim. The gap is the meeting point of formation and selection,
\MT{open:rg-formation-adaptive-family}, with recursion, large networks, nonabelian holonomy, and contact
with any observed system all open.
\end{verdictbox}

\chapter{Related Work}
\label{comp:ch15}

This chapter covers \MT{ch:related-work}, a positioning chapter with one small established remark, and
adds what it does not tabulate: which of the theory's main results are standard mathematics in new
typing and which are new.

\section{The question and the answer}
\label{comp:ch15-question}

The technical chapter asks where the construction sits among existing work and what it owes. The
neighbors flow a distribution along a scale parameter (Bayesian renormalization), coarse-grain graphs
whose edges carry scalars (network renormalization, graph coarsening), constrain learned maps over a base
manifold (equivariant networks), or share, omit, or fit a group-level model (multi-agent active
inference). The present theory computes the parent law from the fine joint by an exact Markov channel at
fixed evidence (\MT{thm:rg-exact-coarse-vfe}, \MT{thm:rg-pointwise-parent-datum}), with group-valued
independent links at one point. The central comparison, with Bayesian renormalization, rests on two
statements: evidence bounds of different joints cannot be compared unless their evidences are related
(\MT{prop:gen-no-distinguished-target}), and an evidence-preserving channel applied to both the posterior
and the recognition law orders the bounds by data processing (\MT{thm:cg-evidence-preserving-channel}).
The scope paragraph states that no other construction is identified with this one and that every
qualification of another work is the technical text's reading.

The link construction adds a direct comparison with lattice gauge theory, induced actions, and
learning rules (\MT{sec:links-related-work}). The references concern particular actions,
normalizations, and updates, rather than an identification of the agent model with a physical
lattice field theory or a neural architecture. The formulas themselves are graded in
\cref{comp:lie-links}.

\section{The picture}
\label{comp:ch15-picture}

\Cref{comp:tab-ch15-map} maps the neighbors; the grade refers to the relation between the neighbor's
central object and the theory's corresponding object.

\begin{sciencetable}{Neighbors of the construction.}
\label{comp:tab-ch15-map}
\small
\begin{tabularx}{\textwidth}{>{\raggedright\arraybackslash}p{0.19\textwidth}>{\raggedright\arraybackslash}X>{\raggedright\arraybackslash}p{0.3\textwidth}>{\raggedright\arraybackslash}p{0.1\textwidth}}
\toprule
\textbf{Neighbor} & \textbf{Its object} & \textbf{Relation here} & \textbf{Grade}\\
\midrule
Bayesian renormalization & Posterior flow on parameter space along a Fisher information scale; stiff and sloppy directions. & Different construction; composition needs a common target (\SG{open:grg-bayesian-bridge}). & loose \\
Information-geometric renormalization & Channels from microscopic to observable, a metric on Hamiltonians, monotone correlation functions. & Only the classical Fisher contraction for one kernel (\MT{thm:cg-fisher-contraction}). & structural \\
Network (Laplacian) renormalization & Heat kernel $e^{-tL}/\Tr e^{-tL}$, spectral entropy, specific-heat peaks. & Spectra gauge invariant, node affinities not (\MT{rem:grg-heat-kernel-gauge}). & structural \\
Algebraic multigrid & Galerkin coarse operator $P^\top AP$ with aggregation prolongation. & Same operator, inside normalized laws (\MT{prop:cg-message-passing-residual}). & exact \\
Matrix-weighted consensus & Laplacian with matrix edge weights; consensus subspace. & The Gaussian interaction family (\SG{def:gauss-interaction-family}). & exact \\
Cellular sheaves & Sheaf Laplacian; $H^0$ as its kernel; Kron obstruction. & Represented-link energies are sheaf energies (\MT{prop:cg-kernel-holonomy}). & exact \\
Connection Laplacians & Orthogonal transports on a graph; convergence to a manifold operator. & Special case: orthogonal links, scalar weights. & exact \\
Maximal-tree gauge fixing & Tree links set to the identity; non-tree links carry loop holonomies. & The rooted datum (\MT{prop:rg-rooted-datum}). & exact \\
Lattice gauge theory & Wilson plaquette term and temporal-gauge time-like loop. & Metric-isometric loop defects match the plaquette term; history loops use identity time-like links (\MT{prop:link-loop-defect}, \MT{prop:link-history-prior}). & exact in scope \\
Induced gauge theory & Integrating out matter produces a link action. & Jointly normalized undirected matter induces an action; record extension does not (\MT{prop:link-induced-action}). & structural \\
Bayesian learning rule & Natural-gradient variational update in Gaussian natural coordinates. & Explicit Euler step of the Gaussian generator Fisher flow (\MT{prop:link-recognition-flow}). & exact \\
Equivariant networks & Weight sharing over a global group or local frames on a base manifold. & Links may be inferred as generators under priors at one point; no learned layer is thereby identified. & loose \\
Multi-agent active inference & Shared models, collective motion, group-level blankets. & Parents computed from the fine joint. & loose \\
Path-space averaging & Relative-entropy bounds for effective diffusions. & Not growing-history, switching-source comparisons. & structural \\
\bottomrule
\end{tabularx}
\end{sciencetable}

The pattern is that the linear algebra is almost entirely standard and the probabilistic and gauge typing
around it is the theory's own. Represented-link quadratic energies are sheaf Laplacian energies with
rectangular restriction maps, whose kernels are global sections; the coarse operator is that of algebraic
multigrid; loops are handled by maximal-tree gauge fixing. The theory places these operators inside
normalized joint laws with exact evidence accounting and two channels that may carry inequivalent
representations, and proves formation theorems there. Its relation to renormalization in the physicist's
sense is weaker than the vocabulary suggests: the main manuscript's renormalization is an exact
finite-step channel with no rescaling, fixed points, or flow of couplings, and the asymptotic apparatus
lives in Supplement~B.

For links, the lineage is more direct but the scope differs. The technical text relates the
metric-isometric loop defect to the Wilson plaquette normalization
\citep{WilsonConfinement1974,Creutz1983}, and the identity time-like links in the history prior to
temporal gauge \citep{KogutSusskind1975,Creutz1983}. The prior uses a Gaussian reference on the Lie
algebra in place of Haar measure; this supplies finite normalization for noncompact groups at the
cost of anchors, covariances, and a fiber metric that transform under reference changes
(\MT{prop:link-generator-prior}). A continuous class function on the general linear group cannot
separate unipotent holonomy from the identity, so a trace-only reading of every loop defect would
miss content that the metric defect detects (\MT{cor:link-no-invariant-defect}).

\begin{analogy}{exact}
Kron reduction of a Gaussian network is integrating out Gaussian sites: the Schur complement of the
precision is the precision of the marginal. The Kron obstruction of \citet{HansenGhrist2019} says that
for vector stalks an eliminated site can induce couplings no pairwise restriction maps express, much as
decimation generates couplings outside the original family. Supplement G classifies when matrix-weighted
Kron reduction stays in its family (\SG{prop:cg-kron-leaves-family}, \SG{thm:cg-congruence-diagonal-kron});
the scalar case is closed (\SG{lem:cg-scalar-kron-closure}). Galerkin closure
(\MT{thm:cg-rectangular-endpoint-closure}) avoids the obstruction by restricting the form to a subspace:
block-spin averaging rather than decimation.
\end{analogy}

\begin{analogy}{exact}
\MT{rem:grg-heat-kernel-gauge} is the lattice-gauge fact that a propagator between different sites is not
gauge invariant. Under vertex reframings $P=\operatorname{diag}(P_i)$ the generator $R^{-1}L$ is
conjugated, so spectrum, heat trace, and spectral entropies are invariant, while off-diagonal heat-kernel
blocks transform as $P_i^{-1}(\cdot)P_j$, like $\langle\phi_i\phi_j^\dagger\rangle$ without a connecting
Wilson line. A partition selector built on node affinities therefore needs invariant block statistics.
\end{analogy}

\section{Objects and notation}
\label{comp:ch15-notation}

\begin{decoder}
Information scale & In Bayesian renormalization, a scale on parameter space set by Fisher distinguishability. & A renormalization scale, by its authors' analogy. \\
Stiff and sloppy directions & Parameter directions well and poorly constrained by data. & Relevant and irrelevant directions, loosely. \\
$D^{\mathfrak O}_\theta$ & Pushed inverse Fisher metric, the diffusion tensor of the data-space flow. & Diffusion constant of a Fokker--Planck equation. \\
$\chi(t)$ & Spectral specific heat $-d\mathcal S/d\log t$; peaks mark resolution scales. & Specific heat. \\
Sheaf, stalk, restriction map & A vector space per vertex and edge, maps from vertex to edge spaces. & Matter field and parallel transport on a lattice. \\
$H^0$ & Global sections: vertex data consistent across every edge. & Covariantly constant configurations. \\
Kron reduction & Eliminating vertices by a Schur complement. & Decimation of Gaussian sites. \\
Galerkin form $P^\top AP$, aggregation & Restricting a form to a subspace of piecewise-constant block fields. & Kadanoff block-spin averaging. \\
Maximal-tree gauge & Gauge in which the links of a spanning tree are the identity. & Axial or tree gauge. \\
Evidence-preserving channel & Coarse channel that leaves the observation marginal unchanged. & Integrating out at fixed partition function. \\
Metric loop defect $\mathfrak Q_p$ & Squared displacement of the represented holonomy from identity in a declared fiber metric. & Wilson plaquette term only for metric isometries. \\
Induced action $-\log\mathsf Z_{\mathrm Y}$ & Contribution obtained by integrating matter out of a jointly normalized undirected model. & Matter-induced gauge action. \\
Generator Fisher flow & Natural-gradient evolution of Gaussian recognition means and precisions. & Variational Bayesian learning rule upon Euler discretization. \\
\end{decoder}

``Renormalization'' in the main manuscript means exact finite-step coarse-graining of the latent law;
Wilsonian words appear only when describing other authors' work. ``Channel'' means a classical Markov
kernel, not the quantum channel of \citet{beny2015information}. ``Gauge'' means passive re-expression of
frames at one point with independent links, not a connection on a base manifold, which is why the
comparison with gauge-equivariant networks \citep{cohen2019gauge} is loose.

\section{The formulas, explained}
\label{comp:ch15-formulas}

The display of Bayesian renormalization the technical text engages is the pushed inverse Fisher metric
\[
D^{\mathfrak O}_\theta=\partial\mathfrak O(\theta)g^F(\theta)^{-1}\partial\mathfrak O(\theta)^\top
\]
(\MTeq{eq:rg-bks-diffusion}), with $\mathfrak O$ the map from parameters to data. In the late-data regime
the posterior is Gaussian with covariance proportional to $(g^F)^{-1}$; the delta method pushes it to data
space with covariance proportional to $D^{\mathfrak O}_\theta$, which \citet{BermanKlingerStapleton2023}
read as the diffusion tensor of a flow on data. The technical text's reading is that this is the leading
local approximation of the pushforward, exact for affine $\mathfrak O$, and that continuing it beyond the
late-data regime is a renormalization construction, not a theorem of Bayesian asymptotics.

The comparison turns on \MT{prop:gen-no-distinguished-target}: for joints indexed by recognition laws
$Q$ and $Q'$,
\[
\Lelbo(Q';X)-\Lelbo(Q;X)
=\mathsf e(Q')-\mathsf e(Q)
-\bigl[\KL(Q'\Vert P_{\theta,Q'}(\cdot\given o,X))-\KL(Q\Vert P_{\theta,Q}(\cdot\given o,X))\bigr],
\]
with $\mathsf e$ the log evidence of each joint. A change of bound mixes a change of evidence with a change
of tightness, so the optimized bounds of two independently declared joints, such as a model at one scale
and a model whose prior is another scale's posterior, are not ordered. An evidence-preserving channel fixes
$\mathsf e$, and data processing (\MT{thm:cg-kl-dpi-extended}) then orders the divergences. In physics
terms, variational free energies of different Hamiltonians cannot be compared unless their partition
functions are related, and integrating out keeps the partition function fixed. Hence the typing
prohibition does not refute Bayesian renormalization, but inserting a posterior as the generative top
factor at another scale declares a new model, comparable only through a supplied comparison map.

For network renormalization the key object is the spectral specific heat
$\chi(t)=-d\mathcal S/d\log t$ of \citet{villegas2023laplacian} (\MTeq{eq:grg-heat-capacity}). Spectral
quantities of the reframed generator survive frame changes, but a scale diagnostic selects no partition
(\MT{open:cg-natural-selector}), and the node affinities one would use to build blocks are not invariant.
For algebraic multigrid the coarse operator $P^\top AP$ with $P$ piecewise constant on disjoint aggregates
\citep{Notay2010Aggregation} is the theory's Galerkin form. The technical documents add a criterion for
exact coarse message passing, invariance of the $M$-orthogonal complement of the coarse range under the
fine operator (\MT{prop:cg-message-passing-residual}), and the placement of the operator inside normalized
laws with gauge covariance and exact evidence accounting.

\section{What is proved, and what is standard}
\label{comp:ch15-scorecard}

The chapter proves one remark; the rest is attribution. The scorecard adds the statements it leans on and
the companion's reading of which main results are standard in substance, which is information, not
criticism: it tells the reader what a result can and cannot do.

\begin{scorecard}
\MT{rem:grg-heat-kernel-gauge} & Spectra and heat traces of the reframed generator are invariant; off-diagonal blocks are not. & \SOLID & Elementary similarity; standard in substance. \\
\MT{prop:gen-no-distinguished-target} & Differences of bounds of different joints mix evidence and tightness. & \SOLID & Algebra of the evidence identity; the point is the typing. \\
\MT{thm:cg-evidence-preserving-channel} & A recognition-independent channel preserves evidence and orders bounds. & \SOLID & Standard in substance (data processing). \\
\MT{thm:elbo-extended-gap} & One evidence bound whose gap is the relative entropy to the posterior, in the extended reals. & \SOLID & Standard in substance (Gibbs variational principle); new in its measure-level, many-agent typing. \\
\MT{thm:rg-exact-coarse-vfe} & Fine free energy equals coarse free energy plus discarded conditional information. & \SOLID & Standard in substance (chain rule). \\
\MT{thm:rg-effective-action} & Coarse likelihood as a conditional partition function; composition along a tower. & \SOLID & Standard in substance (integrating out), with care about versions. \\
\MT{thm:cg-fisher-contraction} & Fisher information contracts under a parameter-independent kernel. & \SOLID & Classical, as the technical text says. \\
\MT{prop:cg-kernel-holonomy} & The kernel of the represented-link energy is the holonomy-fixed sector. & \SOLID & Standard in substance (sheaf $H^0$, cited). \\
\MT{prop:rg-rooted-datum} & Tree links gauge to the identity; non-tree links carry loop holonomies. & \SOLID & Standard (maximal-tree gauge fixing, cited). \\
\SG{prop:gauss-interaction-energy-kernel} & Energy and kernel criterion for matrix-weighted Laplacians. & \SOLID & Standard in substance (consensus literature, cited). \\
\MT{thm:cg-loop-operational-holonomy} & A declared record's law separates loop-holonomy classes. & \SOLID & New as a construction; elementary proof. \\
\MT{prop:link-loop-defect} & Metric-isometric loop defect equals the Wilson plaquette form. & \SOLID & The identification is restricted to isometric represented holonomy. \\
\MT{prop:link-induced-action} & Undirected matter gives an induced link action; record extension leaves the prior generator marginal unchanged. & \SOLID & Normalization distinguishes the two models. \\
\MT{prop:link-recognition-flow} & Gaussian Fisher flow has the Bayesian learning rule as its explicit Euler step. & \SOLID & Stated Gaussian family and step conditions; not a rate-separation theorem. \\
\SG{open:grg-bayesian-bridge} & Relate Bayesian renormalization to Gaussian aggregation through a common monotone. & \OPEN & Needs a common state space and a transported monotone. \\
\end{scorecard}

\begin{hygiene}
The comparisons are claims about other authors' work; their accuracy depends on the cited sources, which
the companion has not rechecked, and the scope paragraph marks each qualification as the technical text's
reading. The Bayesian-renormalization contrast rests only on the evidence identity, not on how the
Gaussian pushforward is read. The sheaf and connection-Laplacian identifications hold for quadratic
energies with represented links, not for the non-Gaussian laws of the general theory, where only the
kernel and the holonomy-fixed sector are claimed.
\end{hygiene}

\section{Standard versus new}
\label{comp:ch15-standard}

Standard in substance, and valuable to know as such, are the evidence-bound identity and its gap (the
Gibbs variational principle), the local-global potential identity (coordinate ascent in mean-field
variational inference, \MT{thm:obs-local-global-potential}), attention as the posterior of a latent
source label (mixture-model responsibilities, \MT{prop:obs-attention-elbo}), the exact coarse free energy
(the chain rule), the effective action (integrating out), Fisher contraction, the kernel of
represented-link energies (sheaf cohomology in degree zero), maximal-tree gauge fixing, the Galerkin
operator of algebraic multigrid, matrix-weighted Laplacians, and the Kron obstruction for vector stalks.
For each, the technical documents supply careful typing: measure-level statements in the extended reals,
two channels with independent links, and a separation of Markov channels, energy restriction, and
parameter aggregation. The typing is the contribution, and it is where the documents' length comes from.

The new latent-link chapter makes further debts explicit (\MT{sec:links-related-work}). Its
Gaussian expectation derivatives use the identities of \citet{OpperArchambeau2009}, and its
Fisher recognition flow gives the Gaussian Bayesian learning rule of \citet{KhanRue2023} as an
Euler step. The comparison-record current has the outer-product form of the fast-weight delta
rule discussed by \citet{Schlag2021}, following \citet{Schmidhuber1992}; the correspondence is
structural because here the current is transported through the derivative of the exponential and
averaged in a latent-generator stationarity condition. Integrating jointly normalized undirected
matter gives an induced action in the sense of the mechanism used by \citet{KazakovMigdal1993}.
The action vanishes for matter attached by normalized record extension; in the Gaussian
connection-Laplacian realization it contains one half a log-determinant plus a self term
(\MT{prop:link-induced-action}). The technical text compares this determinant mechanism with the
cycle-rooted-forest formulas of \citet{Kenyon2011VectorBundleLaplacian}. These are attributions to
specific source formulas, not claims of priority for the full agent construction.

New in content are the formation results: the finite-depth theorem under a supplied hierarchy
(\MT{thm:finite-depth-adaptive-formation}), the endogenous four-cycle witness
(\MT{prop:rg-formation-endogenous-cycle}), the depth-one result for the gated family
(\MT{prop:rg-gated-depth-one}), and the selection and parent line of \cref{comp:ch14}. Their tools are
standard (the Harris theorem, logarithmic Sobolev inequalities, couplings, saddle escape) and their
statements new, within the scopes \cref{comp:ch13,comp:ch14} describe. New as constructions, with
elementary proofs, are the operational holonomy record and the use of the no-distinguished-target identity
against posterior-flow constructions. Supplement G's classification of Kron closure is set against cited
prior art for the obstruction and the scalar case; its assessment belongs to the companion's Supplement G
chapters.

\section{Open ends}
\label{comp:ch15-open}

The bridge between Bayesian renormalization and Gaussian aggregation (\SG{open:grg-bayesian-bridge}) needs
a common state space, a map between the Bayesian scale parameter and a block ratio, and a functional
proved monotone along both constructions, or a Gaussian hierarchy in a stated class with no such
functional. A combined network scheme, sketched but not posed as a numbered problem, would infer a
partition from gauge-invariant diagnostics of edge-event laws through a structural channel with a named
law and push the complete law exactly; its partition part is \MT{open:cg-natural-selector}. The
growing-history, switching-source path comparison that averaging theory does not supply is part of
\MT{open:rg-pointwise-parent-dynamics}.
The finite-dimensional group-valued closure of coarse latent-link laws is
\MT{open:link-coarse-closure}; Gaussian generators and arbitrary real logarithms cannot be
assumed through a general-linear coarse step (\cref{comp:ch12-latent-closure}).

\begin{verdictbox}{\DELIVERED}
The technical chapter says it positions the construction among existing work without identifying it with
any of them, with the Bayesian-renormalization comparison resting on the no-distinguished-target identity
and on evidence-preserving channels. It does that: each comparison is anchored to a statement of the
theory, the Bayesian-renormalization contrast follows correctly from the evidence identity, and the one
proved remark is correct and elementary. Two limits. The chapter does not say, result by result, which of
the theory's statements are standard mathematics in new typing; the companion's reading is that most of
the linear algebra and the variational identities are, and the formation results are new. The added
link comparisons identify the Wilson and temporal-gauge forms at their stated scopes, the
matter-integration mechanism, the delta-rule form, and the Gaussian Bayesian learning update;
they do not establish a physical field or learning architecture. And two
neighbors one would expect here, matrix-weighted consensus and maximal-tree gauge fixing, are treated in
Supplement G and the network-state chapter rather than in this chapter. The attributions to other
authors are the technical text's reading and have not been rechecked by the companion.
\end{verdictbox}
